\documentclass[11pt]{article}
\usepackage[margin=1in]{geometry}
\usepackage{amsmath,amssymb,amsthm}
\usepackage{graphicx}
\usepackage[hidelinks]{hyperref}

\newtheorem{theorem}{Theorem}
\newtheorem{corollary}{Corollary}
\newtheorem{remark}{Remark}

\newcommand{\PhiX}{\Phi_X}
\newcommand{\DeltaPhi}{\Delta_{\Phi}}
\newcommand{\1}{\mathbf 1}

\title{The Indicator Constant on the \(L_p\) Scale}
\author{FA Banach run \texttt{fa\_banach\_001}, agent \texttt{agent\_lane\_11}}
\date{2026-07-03}

\begin{document}
\maketitle

\section*{Status}

This packet records a candidate partial result, likely valid. It answers one
standard family inside item (9) of the ``Looking Forward'' section of
Cwikel--Milman--Rochberg, arXiv:1404.2893. It does not characterize all
K\"othe spaces whose indicator constant is \(\log 2\).

\section*{Source Context}

The source note summarizes Kalton's indicator formalism. For an
indicator-like functional \(\Phi\), it writes
\[
  \Delta_\Phi(f,g)=\Phi(f)+\Phi(g)-\Phi(f+g),
\]
and requires a constant \(\delta(\Phi)\) such that
\[
  0\leq \Delta_\Phi(f,g)
  \leq \delta(\Phi)\bigl(\|f\|_{1}+\|g\|_{1}\bigr).
\]
The note observes that for indicators of K\"othe spaces one always has
\(\delta(\Phi)\leq \log 2\), that equality holds for \(L^1\), and then asks
for what other spaces equality holds.

\begin{center}
\includegraphics[width=0.98\linewidth]{figures/open_problem_crop.png}
\end{center}

\section*{Proof Intuition}

For \(L_p\), the optimizing vector in the indicator definition is forced by a
one-line Lagrange multiplier calculation: \(x^p\) is proportional to \(f\).
Thus the \(L_p\) indicator is exactly \(1/p\) times the \(L_1\) entropy
indicator. The defect \(\Delta_\Phi\) is therefore also scaled by \(1/p\).
Kalton's universal \(\log 2\) bound is just the binary entropy bound, and
disjoint equal-mass functions make this bound sharp.

\section*{Result}

Let \((S,\Sigma,\mu)\) be a measure space. For \(f\geq0\), set
\[
  \Lambda(f)
  = \int_S f\log f\,d\mu-\|f\|_1\log\|f\|_1,
\]
with the usual conventions \(0\log0=0\) and \(\Lambda(0)=0\).

\begin{theorem}[Exact \(L_p\) indicator constants]
\label{thm:lp}
Let \(1\leq p<\infty\). On the natural finite-entropy domain of the
indicator for \(L_p(S,\mu)\),
\[
  \Phi_{L_p}(f)=\frac1p\Lambda(f).
\]
Consequently,
\[
  0\leq \Delta_{\Phi_{L_p}}(f,g)
  \leq \frac{\log 2}{p}\bigl(\|f\|_1+\|g\|_1\bigr).
\]
If the measure algebra contains two disjoint sets of finite positive measure,
then the optimal constant is exactly \((\log 2)/p\). For \(L_\infty\), the
optimal constant is \(0\).
\end{theorem}

\begin{proof}
Fix \(1\leq p<\infty\) and \(f\geq0\), \(f\neq0\), with finite entropy.
Writing \(m=\|f\|_1\), the indicator is
\[
  \Phi_{L_p}(f)=
  \sup_{\|x\|_p\leq1}\int_S f\log|x|\,d\mu.
\]
The strict concavity of \(x\mapsto \log x\) and the constraint
\(\int |x|^p\,d\mu\leq1\) give the maximizer
\[
  x_f=\left(\frac{f}{m}\right)^{1/p}
\]
on the support of \(f\). Therefore
\[
  \Phi_{L_p}(f)
  =\frac1p\int_S f\log\frac{f}{m}\,d\mu
  =\frac1p\Lambda(f).
\]
The case \(f=0\) is immediate.

It remains to recall the sharp entropy-defect bound for \(\Lambda\). Let
\(a=\|f\|_1\), \(b=\|g\|_1\), \(h=f+g\), and \(m=a+b\). On the set where
\(h>0\), put \(u=f/h\), and define the binary entropy
\[
  H(t)=-t\log t-(1-t)\log(1-t).
\]
A direct expansion gives
\[
  \Delta_{\Lambda}(f,g)
  = m H(a/m)-\int_S h\,H(u)\,d\mu .
\]
Since \(H\) is concave and \(\int h u\,d\mu=a\), Jensen's inequality gives
\(\int hH(u)\,d\mu\leq mH(a/m)\). Hence \(\Delta_\Lambda(f,g)\geq0\).
Also \(H(a/m)\leq \log 2\), so
\[
  \Delta_{\Lambda}(f,g)\leq m\log2.
\]
Multiplying by \(1/p\) proves the displayed upper bound.

If \(E,F\) are disjoint measurable sets with \(0<\mu(E),\mu(F)<\infty\), take
\[
  f=\mu(E)^{-1}\1_E,\qquad g=\mu(F)^{-1}\1_F.
\]
Then \(a=b=1\), \(u\in\{0,1\}\) almost everywhere on the support of \(h\), so
\(\int hH(u)=0\) and \(\Delta_\Lambda(f,g)=2\log2\). The ratio
\(\Delta_{\Phi_{L_p}}(f,g)/(\|f\|_1+\|g\|_1)\) is therefore
\((\log2)/p\), proving optimality.

For \(L_\infty\), the constraint \(\|x\|_\infty\leq1\) gives
\(\log |x|\leq0\), while \(x=1\) gives value \(0\). Thus
\(\Phi_{L_\infty}=0\) and the optimal defect constant is \(0\).
\end{proof}

\begin{corollary}[Weighted \(L_p\) spaces]
\label{cor:weighted}
Let \(w>0\) a.e. and \(1\leq p<\infty\). For \(L_p(w\,d\mu)\),
\[
  \Phi_{L_p(w)}(f)
  =\frac1p\Lambda(f)-\frac1p\int_S f\log w\,d\mu
\]
whenever the right-hand side is finite. Hence the same defect bound and the
same sharp constant \((\log2)/p\) hold whenever the measure space admits the
disjoint equal-mass test functions in the indicator domain.
\end{corollary}

\begin{proof}
The constrained maximizer is now
\[
  x_f=\left(\frac{f}{\|f\|_1 w}\right)^{1/p}.
\]
Substitution gives the displayed formula. The additional weight term is
linear in \(f\), so it cancels from
\(\Phi(f)+\Phi(g)-\Phi(f+g)\). The rest is Theorem~\ref{thm:lp}.
\end{proof}

\section*{Implication for the Source Question}

The \(L_p\) scale does not provide new equality cases beyond the \(L^1\)
endpoint. In every nondegenerate measure space covered by the theorem,
\[
  \delta(\Phi_{L_p})=\frac{\log2}{p},
\]
so equality with Kalton's universal upper bound \(\log2\) occurs exactly at
\(p=1\) in this scale. Weighted \(L_1\) spaces have the same extremal constant,
as expected from the cancellation of linear weight terms.

\section*{Limitations and Novelty Check}

This is only a partial result. It does not rule out non-\(L_1\) K\"othe spaces
with \(\delta(\Phi)=\log2\), and it does not address the limiting-case
questions about Kalton's Theorems 1.1 and 6.6 that immediately precede the
indicator-constant question in the source.

Local run indexes were searched for arXiv:1404.2893, Kalton, K\"othe,
centralizer, indicator, \(\delta(\Phi)\), and \(\log2\). No duplicate packet
for this exact calculation was found. A bounded web search for the exact
phrases around ``delta(Phi)'', ``log 2'', ``Kothe'', and ``indicator'' found
the source paper but no exact preexisting answer to this subcase. The
calculation is elementary and may be folklore; the packet records it because it
gives a clean, reviewable partial answer to item (9).

\section*{Human Review Recommendation}

Likely valid. Review should check the conventions in Kalton's indicator
domain and the nondegeneracy assumptions needed for sharpness on the chosen
measure space. The proof of the formula and the upper bound is independent of
those sharpness assumptions.

\begin{thebibliography}{9}

\bibitem{CMR}
M. Cwikel, M. Milman, and R. Rochberg,
\emph{An introduction to Nigel Kalton's work on differentials of complex
interpolation processes for Kothe spaces},
arXiv:1404.2893, 2014.

\bibitem{Kalton}
N. J. Kalton,
\emph{Differentials of complex interpolation processes for Kothe function
spaces},
Trans. Amer. Math. Soc. 333 (1992), no. 2, 479--529.

\end{thebibliography}

\end{document}
